% =========================================================
% Proof Stub: Alternating Series Test
% Source: volume-iii/analysis/sequences/notes/series/notes-series-tests.tex
% Mode: standard
% =========================================================

\clearpage
\phantomsection
\hypertarget{proof-RA-ABB-T-alternating-series}{}

\subsubsection[Alternating Series Test]{Proof --- Alternating Series Test}

\bigskip

\noindent
\textbf{Source.}
Abbott, \emph{Understanding Analysis}; see \texttt{notes-series-tests.tex}.

\vspace{0.75em}

\noindent
\textbf{Goal.}
If $b_n \ge 0$, $(b_n)$ decreasing, and $b_n \to 0$, then $\sum (-1)^{n+1} b_n$ converges.

\vspace{0.75em}

\noindent
\textbf{Logical form.}
\[
  b_n \ge 0,\; b_n \searrow 0 \Rightarrow \sum (-1)^{n+1} b_n \text{ converges}.
\]

\vspace{0.75em}

\noindent
\textbf{Proof strategy.}
Show partial sums $(s_{2n})$ and $(s_{2n-1})$ converge to the same limit. $(s_{2n})$ is increasing and bounded above by $s_1$; $(s_{2n-1})$ is decreasing and bounded below by $s_2$. Use MCT and the fact that $s_{2n-1} - s_{2n} = b_{2n} \to 0$.

\vspace{0.75em}

\noindent
\textbf{Proof.}
\begin{proof}
\Claim If $b_n \ge 0$, $(b_n)$ decreasing, and $b_n \to 0$, then $\sum (-1)^{n+1} b_n$ converges.

\Given 

\Goal 

\Strategy Show partial sums $(s_{2n})$ and $(s_{2n-1})$ converge to the same limit. $(s_{2n})$ is increasing and bounded above by $s_1$; $(s_{2n-1})$ is decreasing and bounded below by $s_2$. Use MCT and the fact that $s_{2n-1} - s_{2n} = b_{2n} \to 0$.

\medskip

% --- s_{2n} = b_1 - b_2 + ... + b_{2n-1} - b_{2n} is increasing (add b_{2n-1} - b_{2n} ≥ 0 each step) ---
% --- s_{2n} ≤ b_1 (bounded above) ---
% --- MCT: s_{2n} → L ---
% --- s_{2n+1} = s_{2n} + b_{2n+1} → L + 0 = L ---
% --- whole sequence s_n → L ---

\end{proof}

\begin{remark*}[Proof shape]
Split into even/odd partial sums; each monotone and bounded; common limit via MCT and $b_n \to 0$.
\end{remark*}

\begin{remark*}[Dependencies]
\begin{itemize}
  \item Monotone Convergence Theorem.
  \item Convergence via Even-Odd subsequences.
  \item $b_{2n} \to 0$ (needed to close the even/odd gap).
\end{itemize}
\end{remark*}
